\subsection{Penance and Reconciliation}

\begin{massbox}{Formal account: judgment placed inside mercy}
Penance is the sacrament in which a baptized sinner, moved by prevenient grace, turns from post-baptismal sin through contrition, confession, and the will to make satisfaction, and receives through a priest's sacramental absolution the forgiveness won by Christ and reconciliation with the Church. Its outward sign is unusual: the penitent's morally sensible acts are its \emph{quasi-matter}, while the absolving word or prayer is form and completion. Christ alone is the principal cause; the priest exercises the keys as an instrumental judge and physician; the penitent does not earn pardon but receives it in a truth-telling conversion that grace itself awakens (John 20:19--23; Trent, sess.~XIV, \emph{De paenitentia}, chs.~1--3; ST III, q.~84, aa.~1--3).
\end{massbox}

\subsubsection{The covenant pedagogy of return}

Sin in Scripture is at once guilt, false worship, broken covenant, enslavement, pollution, wound, debt, and death. No single metaphor exhausts it, and therefore no one-word account exhausts reconciliation. Leviticus joins acknowledgment, restitution, sacrifice, and priestly intercession; after the golden calf Moses intercedes for a people whose communion is ruptured; Nathan names David's concealed act, David confesses, and forgiveness does not prevent temporal consequences (Lev. 5:1--6; 6:1--7; Exod. 32--34; 2~Sam. 12:1--14). Psalm 50 (modern 51) gathers these dimensions into a plea for washing, a created clean heart, restored joy, truthful lips, and the sacrifice of a contrite spirit.

The prophets make clear that return is itself preceded by God. ``Convert us, O Lord, to thee, and we shall be converted'' is not a pious decoration laid over self-reform; it names the first efficient motion of repentance (Lam. 5:21). Jeremiah's new covenant places the law within and promises remembered sin no more; Ezekiel joins sprinkling, a new heart, a new spirit, and the gift of God's Spirit (Jer. 31:31--34; Ezek. 36:25--27). Yet mercy does not render truth unnecessary. The wicked person must turn and live; the injured neighbor must be treated justly; worship cannot substitute for conversion (Ezek. 18:21--32; Isa. 1:10--20).

Christ announces this prophetic imperative as an eschatological gift: repent, because God's reign has drawn near. He forgives the paralytic before the visible healing; receives the woman whose tears, touch, and love manifest faith; restores Zacchaeus into a conversion that includes restitution; and narrates heaven's joy over the found sheep, coin, and son (Mark 1:14--15; 2:1--12; Luke 7:36--50; 15:1--32; 19:1--10). In each case forgiveness is personal without being private. The healed person returns to a people; the lost is carried home; the son is clothed for a household feast.

\begin{threestable}{Paschal moment}{Gift entrusted}{Sacramental consequence}
Keys promised & Peter receives the keys of the kingdom and the authority to bind and loose; the same binding and loosing is addressed to the apostolic community in its work of correction (Matt. 16:18--19; 18:15--20). & Reconciliation possesses an ecclesial judgment; loosing is not an inward feeling severed from the visible Church.\\
Spirit breathed & On Easter evening Christ shows the wounds, gives peace, sends as the Father sent him, breathes the Holy Spirit, and authorizes remission or retention (John 20:19--23). & The sacrament's efficient power is Paschal and pneumatological; retention implies discernment, and discernment ordinarily requires truthful disclosure.\\
Conversion preached & The risen Lord orders repentance for forgiveness to be preached in his name to all nations (Luke 24:46--49; Acts 2:37--42). & The keys serve the universal Gospel. Sacramental confession never replaces faith and conversion; it is their post-baptismal ecclesial form.\\
\end{threestable}

Aquinas therefore locates the sacrament's institution across Christ's mission without fragmenting it. The Lord commands repentance in his preaching, promises the keys, and after the Resurrection manifests both their efficacy and their source in his suffering and rising (ST III, q.~84, a.~7). The absolving sentence does not compete with the Cross. It is a historically extended instrument by which the crucified and risen Christ applies his meritorious cause to this penitent now.

\subsubsection{The Fathers: one mercy through changing disciplines}

The earliest witnesses do not present a modern ritual manual, but they already hold together confession, ecclesial reconciliation, conversion of life, and Eucharistic communion. The \emph{Didache} commands acknowledgment of transgressions ``in the church'' and, for the Lord's Day breaking of bread, confession so that the sacrifice may be pure (4.14; 14.1). These passages directly witness ecclesial confession and Eucharistic preparation; by themselves they do not specify whether every acknowledgment was private, public, sacramental in the later technical sense, or a combination determined by the fault.

The \emph{Shepherd of Hermas}, Mandate IV.3, witnesses both post-baptismal repentance and an austere once-only discipline. Its insistence that the baptized can truly fall and yet receive a divinely provided cure is doctrinally important; its restriction reflects a rigorist disciplinary horizon and is not the Catholic doctrine that later falls are intrinsically unpardonable. The distinction matters: the history of canonical penances includes severe, public, and sometimes unrepeatable forms, while the power of Christ's keys is not exhausted by any one disciplinary form.

Tertullian's Catholic-era \emph{On Repentance} 7--10 makes the medical logic vivid. After baptismal ``shipwreck,'' God leaves an opening for a second repentance; repeated sickness requires repeated medicine; \emph{exomologesis} gives inward repentance an ecclesial body through confession, fasting, tears, prostration, satisfaction, and the intercession of the brethren. His once-only language records early severity, not the final measure of repeatability. His image also guards against cheap grace: the plank is mercy for the shipwrecked, not permission to scuttle the ship. Trent later calls Penance the ``second plank after the shipwreck of lost grace'' while teaching that the sacrament may be received whenever the baptized fall (sess.~XIV, ch.~2 and can.~1).

St.~Cyprian writes amid the Decian persecution, when the reconciliation of the \emph{lapsi} was an urgent ecclesial wound. He urges each to confess while confession, satisfaction, and priestly remission can still be received, and he makes the prophet's ``return with the whole heart'' the measure of the act (\emph{On the Lapsed} 28--29). This is direct pastoral testimony to confession and sacerdotal reconciliation in a public-penitential setting. It does not make the priest the principal cause of pardon; all hope rests on the Lord whose mercy the Church ministers.

The fourth-century conflict with Novatian rigorism clarifies the scope of the keys. St.~Ambrose argues directly from John 20:22--23 that refusing the Church's power to forgive grave sins dishonors the Spirit who gave it; the office is priestly and the power belongs to the true Church (\emph{Concerning Repentance} I.2.5--9). Reading Luke 15, he interprets the returning son's kiss as peace, the robe as the wedding garment, the ring as the Spirit's seal, and the slain calf as Christ's Passion and Eucharistic feast. Thus even grave sinners who truly repent are restored toward sacramental communion (II.3.13--19). This is inherited exegesis, not a claim that each narrative object is a technical component of the rite.

St.~John Chrysostom places the same authority within the fearful dignity of priesthood: priests serve not only at regeneration but afterward, and he adduces James 5:14--15 as forgiveness for the sick after baptism (\emph{On the Priesthood} III.6). St.~Augustine's direct exposition of John 20 says that the Church's Spirit-given love remits the sins of those who participate in it and retains those of persons refusing that communion (\emph{Tractates on John} 121.4); elsewhere he says that the Church received the keys in Peter (124.5). Neither Father makes ecclesial belonging a substitute for conversion. They show why reconciliation with God and reconciliation with Christ's Body cannot be separated.

Finally, St.~Leo the Great's Letter 168.2 (AD 459) forbids the public reading of penitents' written lists and judges it sufficient that troubled consciences be disclosed secretly to priests. The locus is a disciplinary intervention, but it proves that private priestly confession is not a medieval invention and that the mode of disclosure can change while sacramental truth remains. Public penance, private confession, tariffed penances, and today's rites are not identical disciplines; their stable center is conversion submitted to the apostolic keys and ordered to restored communion.

\Needspace{10\baselineskip}
\begin{studybox}{Source boundary: development without anachronism}
The \emph{Didache}, Hermas, Tertullian, and Cyprian directly witness ancient penitential acts and disciplines. Ambrose directly expounds John 20 and Luke 15 against denial of ecclesial forgiveness; Augustine directly expounds John 20 and the Petrine keys; Leo directly regulates secret confession. None should be made to describe every detail of a contemporary confessional. Conversely, later precision about quasi-matter, form, faculty, and integral confession articulates the identity of the sacrament; it does not create Christ's gift centuries after Easter.
\end{studybox}

\subsubsection{Penance as virtue and as sacrament}

The same word names a lifelong virtue and a sacramental act. As virtue, penance is not bare sadness, shame, or self-hatred. Aquinas places it in the will as a species of justice: moved by charity, the sinner detests the past offense against God and purposes amendment and such reparation as friendship permits (ST III, q.~85, aa.~1--4). The past deed cannot become undone; its present disorder can be repudiated, its guilt forgiven, its injuries repaired where possible, and its habits opposed by new acts. Penitential sorrow is therefore rational and hopeful rather than an attempt to reverse time.

As sacrament, this virtuous movement becomes sensible and ecclesial under Christ's keys. Penance's matter is not sin as though evil could be consecrated. Sins are remote matter precisely \emph{as rejected and destroyed}; the proximate quasi-matter is the penitent's contrition, confession, and satisfaction. The priestly act is formal and completive because its words determine these acts toward sacramental absolution (ST III, q.~84, aa.~1--3; q.~90, aa.~1--2). The moral agency of the recipient is thus neither an independent efficient cause of grace nor inert passivity. God moves inwardly; the penitent supplies the acts of a rational subject; Christ through the minister completes the sign.

\begin{fourstable}{Constitutive element}{Metaphysical work}{Moral or ritual act}{Principal guardrail}
Contrition & The will rejects sin as offense and turns toward God under prevenient grace; it supplies the interior root of the quasi-matter. & Sorrow and detestation for sins committed, with a real purpose not to sin and to receive the sacrament when required. & Emotional intensity is not the measure. A purpose of amendment is sincere present resolve, not a guarantee of future impeccability.\\
Confession & Interior judgment becomes a sensible ecclesial accusation submitted to the keys. & After diligent examination, all remembered grave post-baptismal sins are confessed in kind and number, with circumstances that change moral species. & Innocent forgetting does not equal concealment; deliberate withholding of grave sin contradicts integral confession and true conversion. Venial confession is recommended, not strictly required.\\
Satisfaction & Forgiven freedom begins to repair temporal disorder and retrain powers wounded by habit. & The penitent accepts and performs a suitable penance and, independently of the assigned act, owes restitution or repair where justice and prudence require. & Satisfaction has value only in Christ; it neither purchases absolution nor implies that finite acts equal the offense against God. Absolution need not await completion of the assigned penance.\\
Absolution & The priestly word or prayer formally determines the acts toward remission and instrumentally applies the keys. & The approved indicative or deprecatory formula is pronounced by a valid priest intending the Church's act. & God alone forgives principally. Latin analysis of the concluding words ``and I absolve you from your sins in the name of the Father, and of the Son and of the Holy Spirit.'' (CCC 1449) must not invalidate the approved efficacious deprecatory forms of Eastern Churches.\\
\end{fourstable}

Trent distinguishes perfect contrition, arising from charity, from imperfect contrition or attrition, arising from considerations such as sin's ugliness or fear of punishment. Perfect contrition reconciles a person with God before actual reception when it includes the will to receive the sacrament; it does not make the sacrament superfluous. Attrition, when supernatural, hopeful, and joined to a genuine purpose of amendment, is itself God's gift and disposes the sinner to receive justification in the sacrament; fear without conversion, or a plan to continue grave sin, does not (sess.~XIV, ch.~4; CCC 1451--1453).

Aquinas's sign-analysis gives a disciplined account of the sacramental layers. The exterior acts of penitent and absolving priest are the \emph{sacramentum tantum}; inward repentance is his \emph{res et sacramentum}, signified and sacramentally caused or perfected by the complete outward sign without implying that it begins temporally after absolution; forgiveness is the \emph{res tantum} (ST III, q.~84, a.~1 ad 3). Later schools do not allocate the intermediate reality in an entirely uniform way, and Penance imprints no character. The stable point is causal order: outward ecclesial judgment, inward conversion, and divine forgiveness are distinct without becoming three unrelated events.

\subsubsection{Minister, subject, and the stewardship of the keys}

Only a bishop or presbyter validly gives sacramental absolution. A deacon or lay person may hear an anguished disclosure, pray, counsel, and summon a priest; none can complete Penance sacramentally. In Latin discipline, valid absolution ordinarily requires both priestly ordination and the faculty to exercise the keys for this subject (CIC 965--966). This jurisdictional element is fitting to a judicial and ecclesial sign: no minister appoints himself judge of a communion whose keys he received from Christ through the Church.

The capable subject is a living baptized person guilty of post-baptismal sin. Baptism, not Penance, is the unbaptized person's sacrament of first justification. For fruitful and valid reception the penitent must have at least supernatural attrition, intend amendment, and submit the sins to the keys in the manner presently possible. All grave sins remembered after diligent examination belong to integral confession; venial sins may be confessed and are a fruitful matter because repeated sacramental grace heals habits and increases charity (CIC 987--988; Trent, sess.~XIV, chs.~4--5).

\begin{studybox}{Validity, liceity, and emergency discipline}
\textbf{Faculty.} In the Latin Church ordination alone does not ordinarily suffice for valid absolution; the faculty is required. \textbf{Danger of death.} Any priest, even one lacking the ordinary faculty, validly and licitly absolves any penitent from sins and censures in danger of death, even when an approved priest is present (CIC 976). \textbf{General absolution.} It is not a convenience for crowds: Latin law reserves it to imminent danger of death or grave necessity; valid reception requires proper disposition and the intention to confess each grave sin individually in due time (CIC 961--963). \textbf{A narrow invalidity.} Absolution of an accomplice in a sin against the sixth commandment is invalid outside danger of death (CIC 977). These are Latin canonical specifications; Eastern Catholic law and each approved rite must be consulted for its own discipline.
\end{studybox}

The sacramental seal is absolute: a confessor may not betray a penitent in any manner or use confessional knowledge to the penitent's detriment (CIC 983--984). The seal is not a fourth element of matter or form, so its later violation does not retroactively nullify an absolution; it is the grave juridical and moral protection owed to an act in which conscience is opened to Christ's judgment. Likewise, the confessor's personal sanctity is not the efficient source of forgiveness. A sinful but validly ordained and authorized priest can absolve validly because Christ acts; his sin never licenses abuse, doctrinal distortion, imprudent questioning, or violation of the forum's freedom.

The priest is both judge and physician (CIC 978). ``Judge'' protects objective moral truth, integral accusation, the authority to bind or loose, and the need for disposition. ``Physician'' protects the end of judgment: salvation, not humiliation or information-gathering. The two offices correct one another. Mercy without truth abandons the wound; judgment without medicinal charity forgets why Christ entrusted the keys.

\subsubsection{Sacramental operation and the hierarchy of ends}

The causal layers and the hierarchy of ends answer different questions. The sequence from outward sign to \emph{res et sacramentum} to \emph{res tantum} identifies what the sacrament signifies and effects; primary, secondary, and ultimate ends identify why those effects are given and how they are ordered. Thus the primary end is not a fourth reality added after the \emph{res tantum}: it is the remission and restored friendship considered under final causality, as the proper good for which Christ instituted this post-baptismal medicine.

\begin{fourstable}{Ordered level}{Sacramental operation}{Teleological status}{Controlling locus}
Outward sign & Contrition, confession, and satisfaction are the penitent's morally sensible \emph{quasi-matter}; priestly absolution is the form that completes their submission to the keys. & Christ uses the whole ecclesial sign instrumentally toward pardon; the penitent's acts neither merit nor manufacture the grace conferred. & Trent XIV, chs.~2--3; ST III, q.~84, aa.~1--3; q.~90, aa.~1--2.\\
Thomistic \emph{res et sacramentum} & Inward repentance is signified and sacramentally caused or perfected by the complete outward sign as conversion ordered toward remission. It is an intermediate sign-and-reality, not a character. & Converted agency belongs within sacramental reception: it disposes the rational subject under grace without becoming the principal efficient cause of forgiveness. This is causal and logical order, not a claim that contrition begins after absolution. & ST III, q.~84, a.~1 ad 3; later scholastic allocations are not wholly uniform.\\
\emph{Res tantum} & In a disposed penitent, remission of sins and the restoration or increase of sanctifying grace reconcile the sinner with God and restore charity and the virtues flowing with it. & This is the proper sacramental effect considered as reality caused; it is not reducible to legal acquittal or subjective relief. & ST III, q.~84, a.~1 ad 3; q.~89, a.~1; Trent XIV, ch.~3; CCC 1468.\\
Primary proper / proximate end & Reconciliation with God: guilt is remitted and the baptized sinner is restored to grace and charity if these were lost, or strengthened in the filial friendship that remains. & This is the same central gift considered as final cause---the end for which the post-baptismal sign operates. Every further effect flows from or serves it. & Trent XIV, ch.~3; ST III, q.~89, a.~1; CCC 1468, 1496.\\
Subordinate secondary ends & Reconciliation with the Church; peace and serenity of conscience; remission of eternal punishment and at least part of temporal punishment; medicinal strength, repair, and retraining of wounded habits. & These are real sacramental or sacramentally ordered goods, not optional sentiment. They remain subordinate to restored charity: felt peace may lag, ecclesial repair may require action, and temporal consequences may remain. & ST III, q.~86, aa.~4--5; Trent XIV, chs.~3, 8; CCC 1469--1470, 1496.\\
Common ultimate end & Restored Eucharistic communion on pilgrimage and, beyond every sacramental sign, perfected charity in the Father's house and the vision of God. & Penance's proper and secondary ends share the final end of all the sacraments: divine glory and the creature's beatitude in Trinitarian communion. & ST III, q.~60, a.~3; q.~65, a.~3; CCC 1468--1470; Luke 15:20--24; Rev. 21:1--4.\\
\end{fourstable}

The primary effect is reconciliation with God through the remission of sins and restoration or increase of sanctifying grace. Because grace is the created principle from which infused virtues flow, Penance restores the virtues lost with charity. Aquinas carefully distinguishes sacrament and virtue: Penance as sacrament instrumentally causes grace by the keys; contrition is the ultimate disposition of the recipient; Penance as virtue is itself grace's effect (ST III, q.~89, a.~1). This prevents both Pelagian self-repair and a merely forensic account in which the person remains inwardly unchanged.

\enlargethispage{\baselineskip}
Every truly repented sin of a living wayfarer is pardonable because no finite evil exceeds Christ's Passion or God's power to turn the will. Aquinas treats despair of pardon as an injury to divine mercy and to the universal sufficiency of Christ's propitiation (ST III, q.~86, a.~1). Yet forgiveness is not permission to presume. Freedom remains mutable; habits and occasions can remain; restitution and patient ascetic work may still be necessary. The grace restored may be greater, equal, or less in intensity than before the fall according to the recipient's disposition, though it is always genuine grace (q.~89, a.~2).

Forgiveness removes eternal punishment due to mortal sin; it ordinarily leaves some temporal punishment and disorder to be healed. This is not a balance unpaid to an unwilling Christ. Guilt, friendship, consequence, and formed habit are distinguishable. David is forgiven yet undergoes consequences; a thief may be absolved yet still owes restitution; a reconciled addict may still need long retraining. Sacramental satisfaction, voluntary penance, works of mercy, endurance of trials, and indulgences all draw their value from Christ and his communion of saints (ST III, q.~86, aa.~4--5; Trent, sess.~XIV, ch.~8; CCC 1471--1479).

Reconciliation with the Church is a subordinate effect, but not an optional or merely social bonus. Sin wounds the Body; the absolving minister acts in its name; ecclesial charity intercedes; restoration opens the person again to the common sacramental good (1~Cor. 5:1--13; 2~Cor. 2:5--11). The effects therefore include peace and serenity of conscience, spiritual consolation, and increased strength for combat, but subjective relief is not the sacrament's criterion of truth. A penitent may leave consoled or still emotionally burdened; valid absolution rests on Christ's action in the sign, while wise pastoral and psychological care may continue.

\subsubsection{The sacrament among the sacraments and at the last end}

Penance presupposes the baptismal character and restores the grace of baptism without repeating Baptism. Confirmation's character also remains through sin, now summoned back into fruitful witness. Orders supplies the ministerial instrument of the keys; Matrimony and Orders both repeatedly need reconciliation because public vocation does not abolish personal conversion. Anointing may forgive sins when a sick person cannot receive Penance and places no obstacle, but it does not make confession optional when Penance can be received. Above all, Penance tends to the Eucharist: it repairs communion so that the member may receive the sacrament of unity without contradicting the gift (1~Cor. 11:27--32; CCC 1385, 1457).

St.~Ambrose's reading of Luke 15 makes this order visible. The son does not return merely to a clean legal file; he is kissed, clothed, sealed, shod, and brought to the feast whose calf Ambrose relates to the Passion and Eucharist (\emph{Concerning Repentance} II.3.13--19). The sacrament is therefore both retrospective and prospective. It judges and destroys past sin in order to free a future act of charity; it restores the pilgrim to the road and table.

Its anagogical form is equally exact. Sacramental judgment anticipates the final judgment as mercy: the sinner accuses himself now before Christ's minister so as to be loosed, healed, and converted before the will is fixed at death. Satisfaction trains desire for a communion in which no unpaid injury or disordered attachment remains. Absolution gives peace, but not the peace of lowered destiny; it restores the capacity for the highest end. When the Church uses the approved Latin formula reproduced at CCC 1449 or prays an approved Eastern absolving prayer, the final cause is the Father's house: not simply innocence recovered, but filial charity perfected in the banquet where the lost are found and God rejoices in his saints.

\begin{massbox}{The whole motion in one view}
\textbf{The Father} first turns the heart and awaits the child. \textbf{The Son} merits every pardon, institutes the keys, and acts as principal cause through his priest. \textbf{The Spirit} gives contrition, vivifies the ecclesial act, and restores charity. \textbf{The penitent} detests, confesses, accepts repair, and receives rather than manufactures grace. \textbf{The Church} judges medicinally, intercedes, absolves instrumentally, and welcomes to communion. The motion begins in prevenient mercy and ends, if not resisted, in the free friendship of beatitude.
\end{massbox}
